\documentclass[../main.tex]{subfiles}\begin{document}
\section*{User Stories}

\subsection*{Epic 1 -- Konto \& Identitet}

\paragraph{US-01}
Som \textbf{user} vil jeg oprette en konto som \textbf{contributor/requester, organization} eller \textbf{reward partner}, så jeg kan deltage på platformen.

\paragraph{US-02}
Som \textbf{user} vil jeg vedligeholde en profil med relevante oplysninger om mig selv, så andre deltagere kan træffe en informeret beslutning om at interagere med mig.

\subsection*{Epic 2 -- Peer-to-peer-opgaver}

\paragraph{US-03}
Som \textbf{requester} vil jeg oprette en opgave, der beskriver den hjælp, jeg har brug for, så \textbf{contributors} kan tilbyde at udføre den.

\paragraph{US-04}
Som \textbf{contributor} vil jeg gennemse tilgængelige opgaver, så jeg kan finde opgaver, der er relevante for mig.

\paragraph{US-05}
Som \textbf{contributor} vil jeg afgive et tilbud på en opgave, så \textbf{requesteren} kan vælge mig til at udføre den.

\paragraph{US-06}
Som \textbf{requester} vil jeg vælge en \textbf{contributor} blandt de modtagne tilbud, så opgaven kan tildeles.

\paragraph{US-07}
Som \textbf{requester} og \textbf{contributor} vil jeg markere opgaven som fuldført, så transaktionen kan afsluttes.

\paragraph{US-08}
Som \textbf{user} vil jeg anmelde en anden deltager efter en fuldført opgave, så fremtidige brugere kan vurdere vedkommendes omdømme.

\paragraph{US-09}
Som \textbf{user} vil jeg se en anden brugers omdømme, så jeg kan vurdere, om jeg har tillid til vedkommende.

\subsection*{Epic 3 -- Aktiviteter med social impact}

\paragraph{US-10}
Som \textbf{organization} vil jeg oprette frivillige aktiviteter, så \textbf{contributors} kan finde måder at støtte vores sag på.

\paragraph{US-11}
Som \textbf{organization} vil jeg se, hvem der har tilmeldt sig min aktivitet, så jeg kan håndtere deltagelsen.

\paragraph{US-12}
Som \textbf{contributor} vil jeg gennemse tilgængelige frivillige aktiviteter, så jeg kan finde muligheder, der matcher mine interesser og min tilgængelighed.

\paragraph{US-13}
Som \textbf{contributor} vil jeg tilmelde mig en frivillig aktivitet, så \textbf{organizationen} ved, at jeg har til hensigt at deltage.

\paragraph{US-14}
Som \textbf{organization} vil jeg verificere, at en \textbf{contributor} har gennemført en aktivitet, så der kun tildeles Karma Koin for legitime bidrag.

\subsection*{Epic 4 -- Karma Koin}

\paragraph{US-15}
Som \textbf{contributor} vil jeg modtage Karma Koin efter en verificeret aktivitet, så mine sociale bidrag bliver anerkendt.

\paragraph{US-16}
Som \textbf{user} vil jeg se min nuværende Karma Koin-saldo, så jeg ved, hvor mange point jeg kan bruge.

\paragraph{US-17}
Som \textbf{user} vil jeg se min Karma Koin-transaktionshistorik, så jeg kan forstå, hvordan min saldo har ændret sig.

\paragraph{US-18}
Som \textbf{platform} vil jeg have, at enhver Karma Koin-transaktion kan spores tilbage til sin oprindelse, så belønningssystemet kan auditeres.

\subsection*{Epic 5 -- Belønninger}

\paragraph{US-19}
Som \textbf{reward partner} vil jeg oprette belønninger, så \textbf{contributors} kan indløse deres Karma Koin hos min virksomhed.

\paragraph{US-20}
Som \textbf{user} vil jeg gennemse tilgængelige belønninger, så jeg kan beslutte, hvordan jeg vil bruge mine Karma Koin.

\paragraph{US-21}
Som \textbf{user} vil jeg indløse Karma Koin til en belønning, så jeg får en fordel for mit bidrag.

\paragraph{US-22}
Som \textbf{reward partner} vil jeg se statistik for en given belønning, så jeg kan følge, hvor populære belønninger er sammenlignet med hinanden.

\subsection*{Epic 6 -- Tillid \& Sikkerhed}

\paragraph{US-23}
Som \textbf{platform} vil jeg have, at begrænsede handlinger kræver den rette rolle, så \textbf{users} ikke kan udføre handlinger, de ikke er autoriseret til.

\paragraph{US-24}
Som \textbf{administrator} vil jeg undersøge omstridte eller mistænkelige aktiviteter og transaktioner, så misbrug af platformen kan efterforskes.
\end{document}